% Evans second edition, section 1.5, printed pp. 11--12; PDF pages 28--29.
% Chinese: printed pp. 8--9; PDF pages 26--27.
\section{Problems}
\label{sec:1-5}

\subsection*{Source problems (Evans, Chapter 1, Problems 1--5)}
The numbering in this list is the source problem numbering. Solutions in
the next subsection are editorial additions.
\begin{enumerate}
\item Classify each of the partial differential equations in \S1.2 as follows:
\begin{enumerate}[label=(\alph*)]
\item Is the PDE linear, semilinear, quasilinear or fully nonlinear?
\item What is the order of the PDE?
\end{enumerate}
\item Let $k$ be a positive integer. Show that a smooth function defined on
$\R^n$ has in general
\[
 \binom{n+k-1}{k}=\binom{n+k-1}{n-1}
\]
distinct partial derivatives of order $k$.

(Hint: This is the number of ways of inserting $n-1$ dividers $|$ within a
row of $k$ symbols $\circ$: for example,
$\circ\circ\,||\,\circ\circ\circ\,|\,\circ\,|\,\circ\circ\circ\,||\,
\circ\circ\circ\circ\,|$.
Explain why each such pattern corresponds to precisely one of the partial
derivatives of order $k$.)
\end{enumerate}

The next exercises provide some practice with the multiindex notation
introduced in Appendix A.
\begin{enumerate}[resume]
\item Prove the Multinomial Theorem:
\[
 (x_1+\cdots+x_n)^k=\sum_{|\alpha|=k}\binom{|\alpha|}{\alpha}x^\alpha,
\]
where $\binom{|\alpha|}{\alpha}:=|\alpha|!/\alpha!$,
$\alpha!=\alpha_1!\cdots\alpha_n!$, and
$x^\alpha=x_1^{\alpha_1}\cdots x_n^{\alpha_n}$.
The sum is taken over all multiindices $\alpha=(\alpha_1,\ldots,\alpha_n)$
with $|\alpha|=k$.
\item Prove Leibniz's formula:
\[
 D^\alpha(uv)=\sum_{\beta\leq\alpha}\binom{\alpha}{\beta}
 D^\beta uD^{\alpha-\beta}v,
\]
where $u,v:\R^n\to\R$ are smooth,
$\binom{\alpha}{\beta}:=\alpha!/[\beta!(\alpha-\beta)!]$, and
$\beta\leq\alpha$ means $\beta_i\leq\alpha_i$ $(i=1,\ldots,n)$.
\item Assume that $f:\R^n\to\R$ is smooth. Prove
\[
 f(x)=\sum_{|\alpha|\leq k}\frac1{\alpha!}D^\alpha f(0)x^\alpha
 +O(|x|^{k+1})\qquad\text{as }x\to0
\]
for each $k=1,2,\ldots$. This is Taylor's formula in multiindex notation.

(Hint: Fix $x\in\R^n$ and consider the function of one variable
$g(t):=f(tx)$.)
\end{enumerate}

\subsection*{Editorial supplement: solutions to Problems 1--5}
\paragraph{Problem 1: classification.}
The following tables use the most restrictive applicable class, for generic
nonlinearities. The order is the largest total derivative order in space
and time. Coefficients and forcing terms are prescribed, not additional
unknowns, unless explicitly stated otherwise.

\begin{center}\small
\begin{tabular}{@{}llc@{}}
\textbf{Scalar linear equation (\S1.2.1a)}&\textbf{Class}&\textbf{Order}\\\hline
1. Laplace & Linear & 2\\
2. Helmholtz & Linear & 2\\
3. Transport & Linear & 1\\
4. Liouville & Linear & 1\\
5. Heat & Linear & 2\\
6. Schrödinger & Linear & 2\\
7. Kolmogorov & Linear & 2\\
8. Fokker--Planck & Linear & 2\\
9. Wave & Linear & 2\\
10. Klein--Gordon & Linear & 2\\
11. Telegraph & Linear & 2\\
12. General wave & Linear & 2\\
13. Airy & Linear & 3\\
14. Beam & Linear & 4
\end{tabular}
\end{center}

\begin{center}\small
\begin{tabular}{@{}llc@{}}
\textbf{Scalar nonlinear equation (\S1.2.1b)}&\textbf{Class}&\textbf{Order}\\\hline
1. Eikonal & Fully nonlinear & 1\\
2. Nonlinear Poisson & Semilinear & 2\\
3. $p$-Laplacian & Quasilinear & 2\\
4. Minimal surface & Quasilinear & 2\\
5. Monge--Ampère & Fully nonlinear & 2\\
6. Hamilton--Jacobi & Fully nonlinear & 1\\
7. Scalar conservation law & Quasilinear & 1\\
8. Inviscid Burgers & Quasilinear & 1\\
9. Reaction-diffusion & Semilinear & 2\\
10. Porous medium & Quasilinear & 2\\
11. Nonlinear wave & Semilinear & 2\\
12. KdV & Semilinear & 3\\
13. Nonlinear Schrödinger & Semilinear & 2
\end{tabular}
\end{center}
\par\Needspace{65mm}
\PDEAnnotation{先约定参数与分支}{表中按一般非线性列出。
$p=2$ 时 $p$-Laplacian 变为线性；$\gamma=1$ 时多孔介质式变为热方程；
$n=1$ 且 $f$ 已知时 Monge--Ampère 式是 $u_{xx}=f$。
忽略这些例外会把分类误当成方程名称的固定属性。}

The $p$-Laplacian entry assumes $1<p<\infty$, $p\ne2$, using its
classical expansion away from $Du=0$ as in \S\ref{sec:1-2}.
The porous medium entry assumes $\gamma>1$ and uses the expansion where
$u>0$. The Monge--Ampère entry assumes $n\geq2$ and given $f$.
If $H(q,x)$ is affine in $q$, Hamilton--Jacobi is quasilinear (in fact
linear in this particular displayed form). If $f(u)$ is affine, the
nonlinear Poisson, reaction-diffusion and nonlinear wave equations become
linear. For the Schrödinger equations the unknown is complex-valued; one
may equivalently use its real and imaginary parts as a real system.

\begin{center}\small
\begin{tabular}{@{}llc@{}}
\textbf{System (\S1.2.2)}&\textbf{Class}&\textbf{Order}\\\hline
a1. Elastic equilibrium & Linear & 2\\
a2. Elastic evolution & Linear & 2\\
a3. Maxwell & Linear & 1\\
b1. Conservation laws & Quasilinear & 1\\
b2. Reaction-diffusion & Semilinear & 2\\
b3. Incompressible Euler & Quasilinear & 1\\
b4. Incompressible Navier--Stokes & Semilinear & 2
\end{tabular}
\end{center}

For a smooth scalar flux, the chain rule gives
\[
 \operatorname{div}\mathbf F(u)=\sum_i(F^i)'(u)u_{x_i},
\]
hence its quasilinear form. For a system with $m$ components,
$\mathbf F:\R^m\to\R^{m\times n}$ and
$(\operatorname{div}\mathbf F(\mathbf u))^a
=\sum_i\partial_{x_i}F^{ai}(\mathbf u)$.
Affine fluxes give linear equations or systems.

For the incompressible flow entries, the unknowns are $(\mathbf u,p)$.
Navier--Stokes has constant coefficients on the second derivatives of
$\mathbf u$; its nonlinear convection term has order one. This is the
sense of the semilinear classification in the table. The divergence equation
has order one, and there are no second derivatives of $p$ in the displayed
system. It is a system of mixed differential orders, not an unconstrained
uniformly parabolic system for all $n+1$ unknowns. Euler has first-order
derivatives with coefficients depending on $\mathbf u$, hence is quasilinear.

\paragraph{Problem 2: counting independent derivative operators.}
\begin{expandedproof}[editorial solution to source Problem 2]
For a $C^k$ function, the order of differentiation does not matter. Thus a
distinct derivative operator of order $k$ is determined by a multiindex
$\alpha\in\{0,1,2,\ldots\}^n$ satisfying
$\alpha_1+\cdots+\alpha_n=k$. Represent $\alpha_i$ by $\alpha_i$ circles
in the $i$th compartment between $n-1$ dividers. Empty compartments are
allowed: they encode $\alpha_i=0$. This is a bijection.

Among the $k+n-1$ positions, choose the $n-1$ divider positions, giving
$\binom{k+n-1}{n-1}=\binom{k+n-1}{k}$. For $n=1$ there are no dividers
and exactly one derivative, as the same formula predicts. This counts
operators; special functions may have several derivatives with equal values.
\end{expandedproof}

\paragraph{Problem 3: the multinomial theorem.}
\begin{expandedproof}[editorial solution to source Problem 3]
Expand the product of $k$ copies of $(x_1+\cdots+x_n)$ by distributivity.
A term $x^\alpha$ is obtained by selecting $x_i$ from exactly $\alpha_i$
of the factors. First choose $\alpha_1$ of the $k$ factors, then
$\alpha_2$ of the remaining factors, and so on. The number of such choices is
\[
 \binom{k}{\alpha_1}\binom{k-\alpha_1}{\alpha_2}\cdots
 \binom{\alpha_n}{\alpha_n}=\frac{k!}{\alpha_1!\cdots\alpha_n!}.
\]
Grouping all terms with the same multiindex proves the formula. The
identity also holds for $k=0$, with the single multiindex $\alpha=0$.
\end{expandedproof}
\par\Needspace{65mm}
\PDEAnnotation{两类组合系数}{多项式展开中 $\binom{k}{\alpha}=k!/\alpha!$，要求 $|\alpha|=k$。
乘积求导中 $\binom{\alpha}{\beta}=\prod_i\binom{\alpha_i}{\beta_i}$，
逐坐标分配导数。后者一般不等于 $\binom{|\alpha|}{|\beta|}$。}

\paragraph{Problem 4: the multiindex product rule.}
\begin{expandedproof}[editorial solution to source Problem 4]
Apply the one-variable product rule $\alpha_i$ times in coordinate $x_i$,
for each $i=1,\ldots,n$. If exactly $\beta_i$ of these derivatives fall
on $u$, the remaining $\alpha_i-\beta_i$ fall on $v$, and the coefficient
is $\binom{\alpha_i}{\beta_i}$. Since mixed derivatives commute under the
smoothness assumption, the resulting term is
$D^\beta u\,D^{\alpha-\beta}v$, with coefficient
\[
 \prod_{i=1}^n\binom{\alpha_i}{\beta_i}
 =\frac{\alpha!}{\beta!(\alpha-\beta)!}.
\]
Summing over all coordinatewise choices $0\leq\beta_i\leq\alpha_i$
proves the identity. For $\alpha=0$ it reduces to $uv=uv$.
\end{expandedproof}

\paragraph{Problem 5: Taylor's formula with a uniform remainder.}
\begin{expandedproof}[editorial solution to source Problem 5]
It suffices that $f\in C^{k+1}$ on a neighborhood of $0$. Choose $r>0$
so the closed ball $\overline B_r(0)$ lies inside this neighborhood. For
$|x|<r$ set $g_x(t)=f(tx)$, $0\leq t\leq1$. Repeated differentiation
and the multinomial theorem give, for $0\leq j\leq k+1$,
\[
 g_x^{(j)}(t)
 =\left(\sum_{i=1}^n x_i\partial_{x_i}\right)^j f(tx)
 =\sum_{|\alpha|=j}\frac{j!}{\alpha!}x^\alpha D^\alpha f(tx).
\]
The one-variable Taylor formula with integral remainder gives
\begin{align*}
 f(x)&=\sum_{j=0}^k\frac{g_x^{(j)}(0)}{j!}+R_k(x),\\
 R_k(x)&=\frac1{k!}\int_0^1(1-t)^k g_x^{(k+1)}(t)\,dt.
\end{align*}
Substituting the derivative formula converts the finite sum into
\[
 \sum_{|\alpha|\leq k}\frac{D^\alpha f(0)}{\alpha!}x^\alpha.
\]
Continuity and compactness give a finite common bound
\[
 M=\max_{\substack{|\alpha|=k+1\\y\in\overline B_r(0)}}
       |D^\alpha f(y)|.
\]
Using $|x^\alpha|\leq|x|^{k+1}$ and
$\int_0^1(1-t)^k\,dt=1/(k+1)$, we obtain
\[
 |R_k(x)|\leq M\left(\sum_{|\alpha|=k+1}\frac1{\alpha!}\right)
 |x|^{k+1}.
\]
The constant depends only on $n,k$ and the indicated local derivative bound,
not on the direction of $x$. This proves the required $O(|x|^{k+1})$.
\end{expandedproof}
\par\Needspace{65mm}
\PDEAnnotation{余项须对方向一致}{逐条射线使用一元 Taylor 公式还不够：余项常数可能依赖方向。
在固定闭球上统一控制所有 $k+1$ 阶导数，才得到真正的多元
$O(|x|^{k+1})$。仅假设 $C^k$ 通常只能得到 $o(|x|^k)$。}
